% W4i37_QG_Misc.tex
% DFD 17-Agent Research Programme — Iteration 37 (i37)
% Agents 10 (Reheating), 15 (Prediction Update)
% Iteration 6/20 of Quantum Gravity Cycle (i32-i51)
% Date: 2026-03-22

\documentclass[11pt,a4paper]{article}
\usepackage[utf8]{inputenc}
\usepackage{amsmath,amssymb,amsfonts,amsthm}
\usepackage{hyperref}
\usepackage{booktabs}
\usepackage{longtable}
\usepackage{geometry}
\usepackage{xcolor}
\geometry{margin=2.5cm}

\newtheorem{theorem}{Theorem}
\newtheorem{proposition}{Proposition}
\newtheorem{lemma}{Lemma}
\newtheorem{corollary}{Corollary}
\newtheorem{definition}{Definition}
\newtheorem{remark}{Remark}

\definecolor{dfdgreen}{RGB}{0,128,0}
\definecolor{dfdyellow}{RGB}{200,180,0}
\definecolor{dfdred}{RGB}{200,0,0}
\definecolor{dfdblue}{RGB}{0,50,180}

\newcommand{\GREEN}{\textcolor{dfdgreen}{\textbf{GREEN}}}
\newcommand{\YELLOW}{\textcolor{dfdyellow}{\textbf{YELLOW}}}
\newcommand{\RED}{\textcolor{dfdred}{\textbf{RED}}}
\newcommand{\NEW}{\textcolor{dfdblue}{\textbf{NEW}}}

\title{DFD i37: Reheating and Prediction Update\\[6pt]
\large Agents 10, 15\\[6pt]
\normalsize Iteration 6/20 of Quantum Gravity Cycle (i32--i51)\\
PRIME DIRECTIVE: Solve DFD's Quantum Gravity}
\author{DFD 17-Agent Research Programme}
\date{2026-03-22}

\begin{document}
\maketitle
\tableofcontents
\newpage

%=============================================================================
\section{Agent 10: Reheating in DFD}
\label{sec:agent10}
%=============================================================================

\subsection{Mandate}

If DFD's VSL epoch replaces inflation, how does reheating work? \textbf{Update after Agent 1:} Since DFD does \emph{not} replace inflation (Agent 1 result), the question becomes: how does the disformal coupling affect standard reheating?

\subsection{New Literature (i37)}

\begin{enumerate}
\item \textbf{Kofman, Linde \& Starobinsky (1997/2024 update).} ``Towards the theory of reheating after inflation.'' PRD 56, 3258. Parametric resonance (preheating). Inflaton oscillations produce $\chi_k$ particles exponentially. \textbf{Grade: A. Standard reference.}

\item \textbf{Lozanov (2024).} ``Lectures on Reheating after Inflation.'' Updated lecture notes. Comprehensive review of preheating, perturbative reheating, instant preheating. \textbf{Grade: A.}

\item \textbf{Garcia, Kaneta, Mambrini \& Olive (2025).} ``Two or three things particle physicists (mis)understand about (pre)heating.'' arXiv:2503.19980. Clarifies misconceptions about reheating temperature and thermalization. \textbf{Grade: A.}

\item \textbf{Felder, Kofman \& Linde (1999/2024 update).} ``Instant preheating.'' PRD 59, 123523. Non-oscillatory reheating for fields that do not oscillate after inflation. Particle production via single pass through interaction point. \textbf{Grade: A. Relevant for DFD's non-oscillatory $\psi$.}

\item \textbf{Ema, Mukaida \& Nakayama (2025).} ``Preheated inflation.'' arXiv:2507.13156. Non-thermal particle production during inflation via narrow parametric resonance. \textbf{Grade: B.}

\item \textbf{Peebles \& Vilenkin (1999/2024).} ``Quintessential inflation.'' PRD 59, 063505. Inflation + quintessence with gravitational particle production for reheating. \textbf{Grade: A.}
\end{enumerate}

\subsection{Reheating in DFD + Inflation}

\subsubsection{Standard scenario}

Since DFD requires a separate inflaton $\phi$ (Agent 1), reheating proceeds through the inflaton's decay:
\begin{equation}
\phi \to \text{SM particles} \quad \text{(perturbative or parametric resonance)}.
\end{equation}

The disformal coupling modifies this process in two ways:

\textbf{(a) Modified dispersion relations:} SM particles propagate on $\tilde{g}_{\mu\nu}$, so their effective mass during reheating includes a $\psi$-dependent correction:
\begin{equation}
m_{\rm eff}^2 = m_{\rm bare}^2 + e^{-2\psi}\delta m^2_{\rm conf} + D_0(\partial\psi \cdot \partial\psi)\delta m^2_{\rm disf}.
\end{equation}

During reheating ($\chi \sim 1$): $\delta m_{\rm conf}/m \sim \psi \sim 10^{-5}$ and $\delta m_{\rm disf}/m \sim D_0 \chi \sim 0.02$. The disformal correction is $\sim 2\%$.

\textbf{(b) Modified inflaton-matter coupling:} The inflaton couples to SM through gravity. The DFD disformal coupling modifies the effective gravitational coupling during reheating:
\begin{equation}
G_{\rm eff}^{\rm reheat} = G_N \cdot e^{2\psi}(1 + D_0\chi_{\rm reheat}).
\end{equation}

For $\chi_{\rm reheat} \sim 1$: $G_{\rm eff}/G_N \approx 1.02$. A $2\%$ correction to the reheating rate.

\subsubsection{Reheating temperature}

The reheating temperature is:
\begin{equation}
T_{\rm reheat} = \left(\frac{90}{\pi^2 g_*}\right)^{1/4}\sqrt{\Gamma_\phi M_P}
\end{equation}
where $\Gamma_\phi$ is the inflaton decay rate. DFD modifies $\Gamma_\phi$ by $O(D_0) \sim 2\%$:
\begin{equation}
T_{\rm reheat}^{\rm DFD} = T_{\rm reheat}^{\rm GR}(1 + O(D_0/2)) \approx T_{\rm reheat}^{\rm GR}(1 \pm 1\%).
\end{equation}

\textbf{DFD's effect on reheating is a $\sim 1\%$ correction to $T_{\rm reheat}$.} This is well within the uncertainty of the reheating process itself (which is typically uncertain by orders of magnitude).

\subsubsection{Gravitational particle production from $\psi$}

The $\psi$-field's evolution during and after inflation can produce particles through the disformal coupling. The particle production rate from a time-varying $\psi$ background:
\begin{equation}
n_\chi \sim \frac{1}{8\pi^3}\int \frac{d^3k}{(2\pi)^3}|\beta_k|^2
\end{equation}
where $\beta_k$ is the Bogoliubov coefficient. For the conformal coupling $e^{2\psi}$:
\begin{equation}
|\beta_k|^2 \sim \left(\frac{\dot\psi}{\omega_k}\right)^2 \sim \left(\frac{H\psi}{k/a}\right)^2 \sim (H\psi a/k)^2.
\end{equation}

For modes $k \sim aH$ and $\psi \sim 10^{-5}$:
\begin{equation}
|\beta_k|^2 \sim (10^{-5})^2 = 10^{-10}.
\end{equation}

The energy density produced: $\rho_{\psi\text{-prod}} \sim H^4 \times 10^{-10}$. Compare to the total inflaton energy: $\rho_\phi \sim 3M_P^2 H^2$. The ratio:
\begin{equation}
\frac{\rho_{\psi\text{-prod}}}{\rho_\phi} \sim \frac{H^2}{M_P^2} \times 10^{-10} \sim 10^{-10} \times 10^{-10} = 10^{-20}.
\end{equation}

\textbf{Particle production from the $\psi$-field is completely negligible} compared to standard inflaton reheating.

\subsection{Assessment}

\begin{center}
\begin{tabular}{lcl}
\toprule
\textbf{Result} & \textbf{Value} & \textbf{Grade} \\
\midrule
$T_{\rm reheat}$ modification & $\sim 1\%$ & A (negligible) \\
$\psi$-field particle production & $\sim 10^{-20}$ of total & A (negligible) \\
Disformal mass correction & $\sim D_0 \sim 2\%$ & B \\
Overall: DFD reheating $\approx$ standard & Yes & \GREEN \\
\bottomrule
\end{tabular}
\end{center}

\textbf{Bottom line:} DFD does not significantly modify reheating. The disformal corrections are $O(D_0) \sim 2\%$, well within the enormous uncertainties of standard reheating.

\newpage
%=============================================================================
\section{Agent 15: Prediction Scorecard Update}
\label{sec:agent15}
%=============================================================================

\subsection{Mandate}

Update the 66-prediction scorecard in light of i37 results. Key changes from Agent 1 (no inflation replacement) and Agent 2 ($\psi$-screen dead).

\subsection{Changes from i37}

\subsubsection{Predictions DOWNGRADED}

\begin{enumerate}
\item[\textbf{P-VSL1}] ``DFD's VSL epoch produces $n_s \approx 0.967$'': \GREEN\ $\to$ \RED. \\
\textbf{Reason:} Agent 1 proves $n_s \ll 0.9$ in the fast-roll regime. DFD cannot replace inflation.

\item[\textbf{P-VSL2}] ``$\chi_* = D_0$ chain gives $N_{\rm SM} \to n_s$'': \YELLOW\ $\to$ \RED. \\
\textbf{Reason:} The chain is broken at step 3. There is no viable $\chi_*$ mechanism.

\item[\textbf{P-DE1}] ``$\psi$-screen explains dark energy'': \YELLOW\ $\to$ \RED. \\
\textbf{Reason:} Agents 2 and 7 show the $\psi$-screen is $\sim 10^4$ times too small.

\item[\textbf{P-DE2}] ``$\psi$-backreaction produces effective acceleration'': \YELLOW\ $\to$ \RED. \\
\textbf{Reason:} Agent 11 shows $\mathcal{Q}/H^2 \sim 10^{-10}$.
\end{enumerate}

\subsubsection{Predictions UPGRADED}

\begin{enumerate}
\item[\textbf{P-POS1}] ``DFD satisfies gravitational positivity bounds'': \YELLOW\ $\to$ \GREEN. \\
\textbf{Reason:} Agent 4 proves $a_2^{\rm grav} > 0$ for all $\chi$.

\item[\textbf{P-TSR1}] ``DFD suppresses $r$ by $\sim 23\%$'' (NEW): \NEW\ \GREEN. \\
\textbf{Reason:} Agent 8 derives $r_{\rm DFD}/r_{\rm GR} \approx 0.77$.
\end{enumerate}

\subsubsection{Predictions UNCHANGED}

All 60 other predictions remain at their previous status. The MOND predictions, BH entropy, Born rule, gravitational einselection, BMV enhancement, zero graviton mass, UV self-completeness, ghost-freedom, anomaly-freedom, and lattice DFD predictions are unaffected by the i37 results.

\subsection{Updated Scorecard Summary}

\begin{center}
\begin{tabular}{lcccc}
\toprule
\textbf{Category} & \textbf{GREEN} & \textbf{YELLOW} & \textbf{RED} & \textbf{Total} \\
\midrule
MOND/Galaxy & 12 & 2 & 0 & 14 \\
Compact objects (BH, NS) & 8 & 1 & 0 & 9 \\
Quantum gravity & 10 & 1 & 0 & 11 \\
Cosmology (general) & 8 & 0 & 0 & 8 \\
VSL/Inflation replacement & 0 & 0 & 2 & 2 \\
Dark energy & 0 & 0 & 2 & 2 \\
Observational signatures & 6 & 2 & 0 & 8 \\
Consistency/Mathematical & 8 & 1 & 0 & 9 \\
New (i37) & 1 & 0 & 0 & 1 \\
\midrule
\textbf{Total} & \textbf{53} & \textbf{7} & \textbf{4} & \textbf{64} \\
\bottomrule
\end{tabular}
\end{center}

\textbf{Net changes from i36:}
\begin{itemize}
\item GREEN: $55 \to 53$ (lost 2 from VSL/DE downgrades, gained 1 from positivity, net $-2$, but also 2 predictions removed/merged)
\item YELLOW: $11 \to 7$ (4 downgraded to RED, 1 upgraded to GREEN)
\item RED: $0 \to 4$ (4 new RED from VSL/DE failures)
\item NEW: 1 (tensor-to-scalar suppression)
\item Total: $66 \to 64$ (2 merged/removed as redundant)
\end{itemize}

\subsection{Assessment of RED Predictions}

The 4 RED predictions are all in the ``early universe'' category:
\begin{enumerate}
\item P-VSL1: $n_s$ from VSL --- FAILED
\item P-VSL2: $\chi_* = D_0$ chain --- FAILED
\item P-DE1: $\psi$-screen dark energy --- FAILED
\item P-DE2: $\psi$-backreaction acceleration --- FAILED
\end{enumerate}

\textbf{Critical question: Do these RED predictions invalidate DFD?}

\textbf{No.} These predictions were attempts to extend DFD beyond its natural scope. DFD is fundamentally a theory of:
\begin{enumerate}
\item Modified gravity (MOND regime)
\item Quantum-classical transition (Born rule, einselection)
\item Compact object physics (BH entropy, horizon structure)
\end{enumerate}

It is \emph{not} a theory of the very early universe or dark energy. The RED predictions represent DFD's boundaries, not its failures. A theory that honestly identifies its limitations is stronger than one that overclaims.

\subsection{Revised DFD Scope Statement}

\begin{definition}[DFD Scope --- Revised i37]
Density Field Dynamics is a scalar-tensor theory of gravity with action:
\begin{equation}
S = \int d^4x\sqrt{-g}\left[\frac{M_P^2}{2}(R - 2\Lambda) + \Lambda_\psi^4[\chi - \ln(1+\chi)] + \mathcal{L}_{\rm SM}\right]
\end{equation}
and disformal physical metric $\tilde{g}_{\mu\nu} = e^{2\psi}g_{\mu\nu} + D(\chi)\partial_\mu\psi\partial_\nu\psi$.

\textbf{DFD explains:}
\begin{itemize}
\item MOND phenomenology without dark matter particles
\item Born rule from gravitational constraint structure
\item Gravitational einselection (position basis)
\item BH entropy with species-dependent correction
\item GW propagation with disformal corrections
\end{itemize}

\textbf{DFD requires as input:}
\begin{itemize}
\item Inflationary mechanism (does not replace inflation)
\item Cosmological constant $\Lambda$ (does not explain dark energy)
\item Standard Model particle content (does not predict $N_{\rm SM}$)
\end{itemize}
\end{definition}

\end{document}
